\clearpage
\VerdierPhysicalPageAnchor{0198}%
\begingroup\small\setlength{\parskip}{0pt}\setlength{\emergencystretch}{1.5em}
\setlength{\abovedisplayskip}{4pt plus 2pt minus 2pt}%
\setlength{\belowdisplayskip}{4pt plus 2pt minus 2pt}%

\VerdierSemanticAnchor{ch3d-2-3-1-proof-b-conclusion}
于是 $v$ 是满态射，这证明了性质 $(E_3)$，并完成了命题的证明。

\VerdierSemanticAnchor{ch3d-2-3-2}
\medskip\noindent\textbf{命题 2.3.2.}
\textit{设 $M$ 是 $\mathcal A$ 的一个对象集合，满足性质
$(E_1)$、$(E_2)$ 与 $(E_3)$。设
$Y\in\operatorname{Ob}\operatorname{comp}(\mathcal A)$ 是一个复形，使得对每个
整数 $n$ 都有 $H^n(Y)\in M$，并且存在整数 $n_0$，使得对
$n\geq n_0$ 有 $H^n(Y)=0$。则存在复形
$X\in\operatorname{Ob}\operatorname{comp}^{-}(\mathcal A)$，其各分量均属于
$M$，并存在拟同构}
\[
\dot u\colon X\longrightarrow Y.
\]

\VerdierSemanticAnchor{ch3d-2-3-2-proof}
首先，由 (1.2.1)，存在复形
$Y'\in\operatorname{Ob}\operatorname{comp}^{-}(\mathcal A)$ 和拟同构
$Y'\longrightarrow Y$。因此可归结到 $Y$ 是
$\operatorname{comp}^{-}(\mathcal A)$ 的对象这一情形。我们将对 $n$ 归纳构造一列态射

\VerdierSemanticAnchor{ch3d-2-3-2-induction}
\[
u_n\colon X_n\longrightarrow X_{n-1}
\]
使其具有下列性质：

\VerdierSemanticAnchor{ch3d-2-3-2-induction-properties}
\noindent 1) 对充分小的 $n$，$X_n=Y$ 且 $u_n=\operatorname{id}$。

\noindent 2) 复形 $X_n$ 的次数 $i\geq -n$ 的分量属于 $M$。
$X_n$ 的次数 $i\geq -n$ 的微分之核与像均属于 $M^\wedge$（2.1.1）。

\noindent 3) 态射 $u_n$ 是拟同构。它同构于态射
$p((\dot u_n)^{-n},-n)$（2.1.1.3）。

\VerdierSemanticAnchor{ch3d-2-3-2-limits}
完成这一构造后，只需令 $X=\varprojlim_n X_n$，并取典范态射
$\dot u$：
\[
\dot u\colon X=\varprojlim_n X_n\longrightarrow\varinjlim_n X_n=Y.
\]

\VerdierSemanticAnchor{ch3d-2-3-2-induction-step}
对充分小的 $n$，令 $X_n=Y$ 且 $\dot u_n=\operatorname{id}_Y$。
假设对每个整数 $n<q$，已经构造了复形 $X_n$ 与具有上述性质 1)、2) 和 3) 的
态射 $\dot u_n$。复形 $X_{q-1}$ 形如

\VerdierSemanticAnchor{ch3d-2-3-2-xqminus1}
\[
X_{q-1}=\cdots\longrightarrow Y^{-q-2}\longrightarrow Y^{-q-1}
\longrightarrow Z^{-q}\longrightarrow L^{-q+1}\longrightarrow L^{-q+2}
\longrightarrow\cdots,
\]
其中 $Y^i$ 是复形 $Y$ 的分量，而 $L^i$ 属于 $M$。于是有满态射

\VerdierSemanticAnchor{ch3d-2-3-2-quotient}
\[
Z^{-q}\longrightarrow
Z^{-q}/\operatorname{Im}(Y^{-q-1}\longrightarrow Z^{-q})
\longrightarrow0
\]
\endgroup
